\chapter{\texorpdfstring{$\Phi$}{Φ} の類型カタログ}
\label{ch:catalog}
%============================================================

\section{構成方針}

第\ref{sec:dof}節の六自由度から生成される空間のうち、
反復的に現れる点および領域を七族に整理して列挙する。
各類型について、信用ポジション $\kappa$ の符号、決済の性質、代表例を示す。
本カタログは網羅性を主張するものではなく、調査対象の選定に用いる作業用の分類である。

\section{族1：単発交換}

$\delta$ が一回、$\pi$ が一回で対応する形式。

\begin{center}
\begin{tabularx}{\textwidth}{@{}llXX@{}}
\toprule
番号 & 類型 & $\kappa$ / 決済の性質 & 例 \\
\midrule
1-1 & 即時交換 & $\approx 0$ / $\delta$ と同時 & 小売、飲食、券売機 \\
1-2 & 前受単発 & $<0$ / $\delta$ に先行 & 受注生産、結婚式、チケット、CF \\
1-3 & 後払単発 & $>0$ / $\delta$ に後行 & B2B掛売、請求書払い \\
1-4 & 分割後払 & $>0$（長期）/ $\delta$ 一回を $\pi$ 多数に分解 & 割賦、BNPL、住宅ローン \\
\bottomrule
\end{tabularx}
\captionof{table}{族1：単発交換の類型一覧。}
\label{tab:family-1-catalog}
\end{center}

1-4 のみ構造が質的に異なる。$\delta$ が一時点であるのに $\pi$ を長期に展開しており、
第\ref{sec:objective}章の割引因子の差を単独で商品化した形式である。
余剰源は $\phi_{\mathrm{barg}}$ と $\phi_{\mathrm{cog}}$ に偏る。

\section{族2：継続アクセス}

権利 $\dbar$ を販売し、$\delta$ が行使される形式。認知余剰の主要な所在。

\begin{center}
\begin{tabularx}{\textwidth}{@{}llXX@{}}
\toprule
番号 & 類型 & $\kappa$ / 決済の性質 & 例 \\
\midrule
2-1 & 定額アクセス & 年払 $<0$ / 月末払 $>0$、$\pi=c$ & SaaS、会員制小売、施設利用型 \\
2-2 & 従量 & $\approx 0$ / $\pi = p\,\delta$ & 電気、クラウド、通信 \\
2-3 & 二部料金 & 混合 / $\pi = F + p\,\delta$ & 携帯、会費＋商品、予約割引 \\
2-4 & プリペイド & $<0$ / $\pi$ 先行、$\delta$ の時期は顧客裁量 & ギフトカード、回数券、交通系IC \\
2-5 & オプション・保証 & $<0$ / $\pi$ 先行、$\delta$ は行使時のみ & 延長保証、キャンセル可予約、融資枠 \\
\bottomrule
\end{tabularx}
\captionof{table}{族2：継続アクセスの類型一覧。}
\label{tab:family-2-catalog}
\end{center}

2-1 は容量制約の有無により二分される（第\ref{sec:breakage}節の命題）。
2-4 と 2-5 は形式的にオプションの売却であり、事業者は売り持ちの状態にある。

\section{族3：資産の時間貸し}

\begin{center}
\begin{tabularx}{\textwidth}{@{}llXX@{}}
\toprule
番号 & 類型 & $\kappa$ / 決済の性質 & 例 \\
\midrule
3-1 & リース・レンタル & 資産が先行 / 時間比例 & 機械リース、不動産賃貸 \\
3-2 & 共用資産 & 資産が先行 / 稼働に連動 & 宿泊、カーシェア、コワーキング \\
3-3 & 譲渡＋長期回収 & $>0$（大）/ 所有権移転後も $\pi$ が継続 & 販売金融、設備ファイナンス \\
\bottomrule
\end{tabularx}
\captionof{table}{族3：資産の時間貸しの類型一覧。}
\label{tab:family-3-catalog}
\end{center}

この族の特徴は $\kappa$ ではなく式\eqref{eq:cash}の $A$ が支配的になる点にある。
3-2 は容量制約が明示的に存在するため、稼働率が実行可能性を決定する。

\section{族4：補完財による回収}

$\delta$ を分割し、$\pi$ を意図的に偏らせる形式。

\begin{center}
\begin{tabularx}{\textwidth}{@{}llXX@{}}
\toprule
番号 & 類型 & $\kappa$ / 決済の性質 & 例 \\
\midrule
4-1 & ロックイン消耗品 & 初期 $>0$ / $\delta_1$ 原価割れ、$\delta_2$ で回収 & 印刷機器、抽出機器、家庭用ゲーム機 \\
4-2 & フリーミアム & 初期 $>0$ / 一部に $\pi=0$、転換分から回収 & ストレージ、業務チャット、モバイルゲーム \\
4-3 & 本体補填＋金融 & $>0$ / 機器を割引し継続料で回収 & 通信キャリア、自動車販売金融 \\
\bottomrule
\end{tabularx}
\captionof{table}{族4：補完財による回収の類型一覧。}
\label{tab:family-4-catalog}
\end{center}

この族は定義上、初期に $\kappa>0$ を意図的に作り出す。
顧客獲得費用の回収期間が長期化すると、族全体が資本市場に依存することになる。

\section{族5：成果・状態連動}

$\pi$ が $\omega$ に可測な形式。第\ref{sec:info}章の議論が直接効く。

\begin{center}
\begin{tabularx}{\textwidth}{@{}llXX@{}}
\toprule
番号 & 類型 & $\kappa$ / 決済の性質 & 例 \\
\midrule
5-1 & 成果報酬 & $>0$（大）/ $\pi=h(\text{結果})$ & 人材紹介、M\&A助言 \\
5-2 & レベニューシェア & $>0$ / $\pi = $ 顧客売上 $\times$ 率 & フランチャイズ、アプリ流通、代理店 \\
5-3 & コストプラス & 中間 / $\pi = $ 実費 $+$ マージン & 一部の建設、コンサル、公共調達 \\
5-4 & 保険引受 & $<0$ / $\pi$ 先取り、$\delta$ が $\omega$ 依存 & 生損保、保証 \\
5-5 & IPライセンス & 混合 / ロイヤルティ、$\delta$ の限界費用ゼロ & 半導体IP、キャラクター許諾、特許 \\
\bottomrule
\end{tabularx}
\captionof{table}{族5：成果・状態連動の類型一覧。}
\label{tab:family-5-catalog}
\end{center}

5-1 と 5-2 は努力が観測できないことによるリスク移転、
5-3 は結果が不確実であることによる逆向きの移転である（第\ref{sec:info}章の注）。
5-4 のみ族内で時間構造が鏡像になっている。

\section{族6：媒介}

第三者間の取引に介在する形式。

\begin{center}
\begin{tabularx}{\textwidth}{@{}llXX@{}}
\toprule
番号 & 類型 & $\kappa$ / 決済の性質 & 例 \\
\midrule
6-1 & 取引額手数料 & $\approx 0$ / 流通額 $\times$ 率 & 証券、不動産仲介、マーケットプレイス \\
6-2 & エスクロー預り & $<0$（他人の資金）/ 一時預かり後に送金 & 決済代行、宿泊仲介、旅行 \\
6-3 & スプレッド & $>0$（在庫）/ 買値と売値の差 & 卸、両替、マーケットメイク \\
6-4 & 機会課金 & $<0$ / 取引ではなく掲載・接触に課金 & 求人広告、検索連動広告 \\
\bottomrule
\end{tabularx}
\captionof{table}{族6：媒介の類型一覧。}
\label{tab:family-6-catalog}
\end{center}

6-2 の $\kappa$ は自己資金ではない。式\eqref{eq:cash}上は他の負の $\kappa$ と同位置に現れるが、
流量減少時に急速に流出する点で性質が異なる。
6-1 と 6-3 の差は在庫リスクを保有するかどうかにあり、これも第\ref{sec:info}章の移転の問題である。

\section{族7：受益者と支払者の分離}

添字 $i$ を動かす操作そのもの。

\begin{center}
\begin{tabularx}{\textwidth}{@{}llXX@{}}
\toprule
番号 & 類型 & $\kappa$ / 決済の性質 & 例 \\
\midrule
7-1 & 広告 & 混合 / 受益者 $\pi=0$、第三者が支払う & 放送、検索、SNS \\
7-2 & クロスサブシディ & 混合 / 片側無料、反対側から回収 & 両面市場の片面課金 \\
7-3 & 公的支払 & $>0$ / 保険者・政府が支払う & 医療、介護、教育 \\
7-4 & 寄付・支援 & $<0$ / 対価関係が明示されない & 支援型課金、投げ銭、非営利 \\
\bottomrule
\end{tabularx}
\captionof{table}{族7：受益者と支払者の分離の類型一覧。}
\label{tab:family-7-catalog}
\end{center}

7-3 は特殊であり、支払者が価格を規定するため $\Phi$ の設計自由度が制度層に吸収されている。
公定価格の改定により $\kappa$ と $m$ が同時に外生的に変化する。

\section{カタログの限界}

以上二七類型は網羅的ではない。当初の想定では三つの欠落があったが、
うち二つは既存類型で記述できる。

\subsection{労働側の \texorpdfstring{$\Phi$}{Phi}：族を追加しない}
\label{sec:labor-phi}

雇用契約もまた履行と決済の写像である。ただし方向が逆であり、
労働者が $\delta$ を提供し、企業が $\pi$ を支払う。
第\ref{sec:phi}章の $X$ の符号規約は企業への流入を正とするため、
カタログの二七類型はすべて企業が $\delta$ を提供する側を想定している。

しかし\textbf{視点を労働者側に置けば、労働者は $\delta$ を提供する事業者である}。
第\ref{ch:solo}章で扱う一人事業と同じ構造であり、
カタログはそのまま適用できる。

\begin{center}
\begin{tabularx}{\textwidth}{@{}lXl@{}}
\toprule
契約形態 & 該当する類型 & $\kappa_L$ \\
\midrule
月給・後払い & 1-3 後払単発の反復 & $>0$（労働者が与信） \\
日払い・週払い & 1-1 即時交換 & $\approx 0$ \\
前借り・支度金 & 1-2 前受単発 & $<0$ \\
賞与 & 5-1 成果報酬 & $>0$ \\
持分報酬 & 5-2 レベニューシェア & $>0$（極大） \\
退職金 & 1-4 分割後払 & $>0$（超長期） \\
\bottomrule
\end{tabularx}
\captionof{table}{労働者側から見た契約形態と該当類型の対応。}
\label{tab:labor-phi-mapping}
\end{center}

\textbf{新たな族は不要である。} 労働者は族1と族5を組み合わせた複合契約を結んでいる。

\begin{remark}[「特殊例」ではなく制約の異なる対象である]
\label{rem:worker-constraints}
表\ref{tab:labor-phi-mapping}は視点を反転すればカタログが適用できることを示すが、
労働者を一人事業の特殊例と見なすことはできない。
$\Phi$ と $\phi$ の両方で制約が異なる。

\textbf{$\Phi$ の制約。} 第\ref{sec:dof}節の六自由度のうち五つが法により制約される。
労働基準法は賃金について通貨払い、直接払い、全額払い、
毎月一回以上、一定期日払いを定めており、
時間関係は後払いに、支払主体は使用者に固定される。
状態依存性も、賃金を業績に連動させることは限定的にしか認められない。
\textbf{選択できないのではなく、選択肢の集合が小さい。}

\textbf{$\phi$ の制約。} 労働者について $\phi_{\mathrm{prod}}$ を定義するには
限界費用が必要だが、これが定まらない。
$\phi_{\mathrm{cog}}$ に至っては対応物が不明である。
第\ref{sec:surplus}章の三分解は労働者にそのまま適用できない。

さらに第\ref{sec:objective}章の倒産時刻 $\tau$ は吸収状態であったが、
労働者は支払不能により消滅しない。
本稿は労働者を独立した対象として扱わない。
\end{remark}

二点を記録しておく。第一に、月給後払いが標準であるため
$\kappa_L>0$ が支配的である。給与の支払サイトは通常15--45日であり、
仕入債務と同じ桁の与信を労働者が企業に供与している。
式\eqref{eq:cash}の第三項に $\kappa_L$ が含まれることを明示すべきである。

第二に、雇用契約は信用の帰属を組織に置く代わりに $\kappa$ の安定を買う取引でもある
（第\ref{sec:credit-formation}節）。これはカタログの言葉では次のように書ける。
\begin{quote}
労働者は 5-1（成果報酬、高分散）を避けて 1-3（後払単発、低分散）を選び、
対価として公開型信用の蓄積を放棄している。
\end{quote}
これは命題\ref{prop:discount}と同型であり、
安定を選ぶ者は将来の請求権を割り引いて手放している。

\subsection{複合契約：既存類型の合成}

B2B の実際の契約は複数類型の合成であることが多い。
これは\textbf{カタログの限界ではなく、正しい使い方である}。
単一の類型に還元できないことは、類型が基底として機能していることを意味する。

\subsection{請求権の流通：枠組みの外にある}
\label{sec:token-limit}

トークン発行のように、$\pi$ が将来の $\delta$ に対する請求権として
それ自体流通する形式は、依然として扱えない。

式\eqref{eq:phi}の $\Phi:\Delta\to\Pi$ は、
履行と決済が同一の相手方との間で対応することを前提としている。
請求権が第三者に転売されると、\textbf{履行を受ける主体が事後的に変わる}。
これは定義域の問題であり、$N$ を固定した第\ref{sec:phi}章の設定では表現できない。
枠組みの拡張を要する。

\section{現実企業は \texorpdfstring{$\Phi$}{Φ} の束である}

規模の大きい企業は通常、複数の族にまたがる $\Phi$ を並行して運用する。
$\kappa$ の符号が族ごとに異なるため、企業単位で $\CCC$ を計算すると
異符号の混合平均を観測することになる。
第\ref{sec:unit}節の結論を再掲すれば、\textbf{観測単位は企業ではなく $\Phi$ に置くべきである}。

%============================================================
